\section{Positive response modes and finite rank}\label{sec:positive-modes}
The single-power test extends to a declared positive finite response class. All following Hankel identities assume point samples or common translated bin windows. Unequal windows require Appendix~\ref{app:statistics}'s window-aware treatment; one true mode can otherwise fail the single-mode identity.
On an equally spaced logarithmic grid $s_j=s_0+jh$, suppose
\begin{equation}
 Y_j=\alpha+\sum_{a=1}^r B_a\lambda_a^j,
 \qquad B_a>0,\quad 0<\lambda_a<1,
 \label{eq:positive-modes}
\end{equation}
with distinct nodes. Here $\lambda_a=e^{-p_a h}$ and $B_a$ includes the
initial-grid factor. These are inference assumptions; their identification
with a spectral measure of the physical source remains a separate derivation.
For $e_j=Y_j-Y_{j+1}$ and $w_a=B_a(1-\lambda_a)$,
\begin{equation}
 \begin{gathered}
 (H_m)_{ij}=e_{i+j},\qquad (V_m)_{ia}=\lambda_a^i,\\
 H_m=V_m\operatorname{diag}(w_a)V_m^T.
 \end{gathered}
 \label{eq:positive-hankel}
\end{equation}
Positivity and the Vandermonde minors give
$H_m\succeq0$ and $\operatorname{rank}H_m=\min(m,r)$.
A positive $r$-square determinant therefore detects at least $r$ visible
modes. An increasing branch has the same result after reversing all
differences when its amplitudes have a common opposite sign. Mixed signed
amplitudes require another analysis.

Write $p(x)=\prod_a(x-\lambda_a)=x^r+\sum_{k<r}c_kx^k$.
The moment recurrence and reconstruction system are
\begin{equation}
 \begin{gathered}
 e_{j+r}+\sum_{k<r}c_ke_{j+k}=0,\\
 H_r\bm c=-(e_r,\ldots,e_{2r-1})^T.
 \end{gathered}
 \label{eq:prony-system}
\end{equation}
Thus $2r+1$ exact $Y$ samples reconstruct $p$, the distinct nodes, the
positive weights through the Vandermonde system, $B_a=w_a/(1-\lambda_a)$,
and $\alpha=Y_0-\sum B_a$. Membership checks on roots and weights are part
of the reconstruction. One additional sample checks the next recurrence.
For a general positive moment measure on $[0,1)$, exact flat rank has the
stronger certificate
\begin{equation}
 \begin{gathered}
 (c_0,\ldots,c_{r-1},1)H_{r+1}(c_0,\ldots,c_{r-1},1)^T\\
 =\int p(\lambda)^2\,\mathrm d\mu=0.
 \end{gathered}
 \label{eq:flat-positive-rank}
\end{equation}
The nonnegative integrand forces support onto the roots; rank $H_r=r$
then gives exactly $r$ visible atoms within this exact positive class.

\subsection{Approximate plateau mass needs a rate restriction}
If a specified polynomial satisfies $\int p^2\mathrm d\mu\le\eta$ and
$|p|\ge d>0$ outside proposed node neighborhoods $S$, then
\begin{equation}
 \begin{gathered}
 \mu(S^c)\le\eta/d^2,\\
 B(S^c)\le\frac{\eta}{d^2(1-\lambda_{\max})}
 \quad\text{if }\lambda\le\lambda_{\max}<1.
 \end{gathered}
 \label{eq:approx-positive-rank}
\end{equation}
The first inequality integrates a nonnegative lower bound; the second uses
$B=\int(1-\lambda)^{-1}\mathrm d\mu$. Without the independently justified
upper node bound, tiny difference weight can conceal finite plateau mass
near $\lambda=1$. This is the same slow-mode obstruction as the finite-window
endpoint construction. A noisy small determinant alone supplies neither exact
closure nor a physical rate cutoff. Adjacent differences reuse observations;
their full covariance and any data-dependent polynomial selection enter a
simultaneous confidence calculation.

\subsection{A compact exact certificate}
A declared dimensionless example uses
$\alpha=33/40$, $B=(1/10,1/20,1/40)$ and
$\lambda=(1/2,1/3,1/4)$. It has $Y_0=1$,
$\operatorname{rank}H_3=\operatorname{rank}H_4=3$, and
\[
 p(x)=x^3-\frac{13}{12}x^2+\frac38x-\frac1{24}.
\]
Eight exact samples retain the seventh-sample reconstruction and one
additional consistency check. The rational checker authenticates every
moment, recurrence and recovered amplitude. Native rank and solve receipts
check the same frozen operands separately. The example is a finite inference
certificate; it does not replace the actual spin-generator closure or supply
cosmological rates from the source ledger.
